\SourcePage{567}{570}
\section{Current density in a magnetic field}

We derive the quantum-mechanical expression for the current density.
The charged particle moves in a magnetic field.
Our starting point is the formula\footnote{In this section, $\mathbf j$ denotes
the electric-current density. It equals the particle-current density multiplied
by the particle charge $e$.}:

\begin{equation}
 \delta H=-\frac1c\int\mathbf j\,\delta\mathbf A\,dV,
 \tag{115.1}
\end{equation}
For charges distributed in space, equation (115.1) determines the change in the
Hamilton function.  The change in question results from varying the vector
potential\footnote{Consider a charge in a magnetic field.  Its Lagrangian
contains the term $(e/c)\mathbf v\mathbf A$.  When the charge is represented as
distributed in space, this term is instead written
$(1/c)\int\mathbf j\mathbf A\,dV$.  Consequently, variation of $\mathbf A$
changes the Lagrangian by
\[
 \delta L=\frac1c\int\mathbf j\,\delta\mathbf A\,dV.
\]
An infinitesimal change in the Hamilton function is the negative of the change
in the Lagrangian.  See Vol.~I, \S~40.}.  In quantum mechanics, it must be
applied to the mean value of the Hamiltonian of a charged particle:

\begin{equation}
 \overline H=\int\Psi^*\left[
 \frac1{2m}\left(\hat{\mathbf p}-\frac ec\mathbf A\right)^2
 -\frac\mu s\mathbf H\hat{\mathbf s}\right]\Psi\,dV.
 \tag{115.2}
\end{equation}
Varying this expression and recalling that
$\delta\mathbf H=\Curl{\delta\mathbf A}$, we obtain
\begin{align}
 \delta\overline H={}&\int\Psi^*\left[-\frac e{2mc}
 (\hat{\mathbf p}\,\delta\mathbf A+\delta\mathbf A\hat{\mathbf p})
 +\frac{e^2}{mc^2}\mathbf A\,\delta\mathbf A\right]\Psi\,dV\notag\\
 &-\frac\mu s\int\Curl{\delta\mathbf A}\cdot
 \Psi^*\hat{\mathbf s}\Psi\,dV.
 \tag{115.3}
\end{align}
We transform the term containing $\hat{\mathbf p}\,\delta\mathbf A$ by
integration by parts:
\[
 \int\Psi^*\hat{\mathbf p}\,\delta\mathbf A\Psi\,dV
 =-i\hbar\int\Psi^*\nabla(\delta\mathbf A\Psi)\,dV
 =i\hbar\int\delta\mathbf A\Psi\nabla\Psi^*\,dV
\]
(as usual, the integral over the surface at infinity vanishes). We also
integrate the last term in (115.3) by parts, using the familiar vector-analysis
identity
\[
 \mathbf a\Curl{\mathbf b}
 =-\Div{(\mathbf a\times\mathbf b)}
 +\mathbf b\Curl{\mathbf a}.
\]

\SourcePage{568}{571}
The integral of the divergence term vanishes, leaving
\[
 \int\Psi^*\hat{\mathbf s}\Psi\Curl{\delta\mathbf A}\,dV
 =\int\delta\mathbf A\Curl{(\Psi^*\hat{\mathbf s}\Psi)}\,dV.
\]
The final result is therefore
\begin{align*}
 \delta\overline H={}&-\frac{ie\hbar}{2mc}\int\delta\mathbf A
 (\Psi\nabla\Psi^*-\Psi^*\nabla\Psi)\,dV\\
 &+\frac{e^2}{mc^2}\int\mathbf A\,\delta\mathbf A\Psi\Psi^*\,dV
 -\frac\mu s\int\delta\mathbf A\Curl{(\Psi^*\hat{\mathbf s}\Psi)}\,dV.
\end{align*}
Comparison with (115.1) gives the following expression for the current
density:
\begin{equation}
 \mathbf j=\frac{ie\hbar}{2m}
 [ (\nabla\Psi^*)\Psi-\Psi^*\nabla\Psi]
 -\frac{e^2}{mc}\mathbf A\Psi^*\Psi
 +\frac\mu s c\Curl{(\Psi^*\hat{\mathbf s}\Psi)}.
 \tag{115.4}
\end{equation}
We emphasize that although the vector potential occurs explicitly in this
expression, the result is nevertheless unambiguous, as it must be. This can be
verified directly by observing that the gauge transformation of the vector
potential must, in accordance with (111.8), be accompanied by the wave-function
transformation (111.9).

It is also readily checked that the current (115.4), together with the charge
density $\rho=e|\Psi|^2$, obeys the continuity equation, as required:
\[
 \frac{\partial\rho}{\partial t}+\Div{\mathbf j}=0.
\]
The final term of (115.4) is the contribution to the current density arising
from the particle's magnetic moment. It has the form $c\Curl{\mathbf m}$, where
\begin{equation}
 \mathbf m=\frac\mu s\Psi^*\hat{\mathbf s}\Psi
 =\Psi^*\hat{\boldsymbol\mu}\Psi
 \tag{115.5}
\end{equation}
is the spatial density of magnetic moment.

Expression (115.4) is the mean value of the current density. It may be regarded
as a diagonal matrix element of an operator, the current-density operator
$\hat{\mathbf j}$. This operator has its simplest form in the
second-quantization representation: one merely replaces $\Psi$ and $\Psi^*$ by
the operators $\hat\Psi$ and $\hat\Psi^+$ (and, by the general rule,
$\hat\Psi^+$ must stand to the left of $\hat\Psi$ in every term). The
off-diagonal matrix elements of this operator can likewise be defined:
\begin{equation}
 \mathbf j_{nm}=\frac{ie\hbar}{2m}
 [ (\nabla\Psi_n^*)\Psi_m-\Psi_n^*\nabla\Psi_m]
 -\frac{e^2}{mc}\mathbf A\Psi_n^*\Psi_m
 +\frac\mu s c\Curl{(\Psi_n^*\hat{\mathbf s}\Psi_m)}.
 \tag{115.6}
\end{equation}
